\section{\texorpdfstring{The high-\(\beta\) ledger and the \(1/400\)
gate}{The high-beta ledger and the 1/400 gate}}
\label{sec:tpc32-high-beta-ledger}

We now specialize the preceding unconditional and conditional statements
to the high-row-scale packet inherited from TPC-28--TPC-31.  Its exponents
are
\begin{equation}
 \boxed{
 \beta=\frac{267}{400},
 \qquad
 1-\beta=\frac{133}{400}.}
 \label{eq:tpc32-high-beta-scales}
\end{equation}
Thus
\begin{equation}
 Q=X^{267/400+o(1)},
 \qquad
 J=X^{133/400+o(1)},
 \qquad
 N_0=JQ^2\asymp XQ.
 \label{eq:tpc32-high-beta-QJ}
\end{equation}
All endpoint statements below retain the usual strict-margin
convention: every fixed exponent below the displayed endpoint is
proved, but equality at an endpoint is not asserted.

\subsection{Exact fraction arithmetic}

The small numerical gap driving this section is best kept in exact
fractions:
\begin{align}
 1-\beta
 &=\frac{133}{400}=\frac{266}{800},
 \label{eq:tpc32-high-large-content-margin}\\
 \frac\beta2
 &=\frac{267}{800},
 \label{eq:tpc32-high-rms-margin}\\
 \frac\beta2-(1-\beta)
 &=\frac{267}{800}-\frac{266}{800}
 =\boxed{\frac1{800}},
 \label{eq:tpc32-high-amplitude-slack}\\
 3\beta-2
 &=\frac{801}{400}-\frac{800}{400}
 =\boxed{\frac1{400}},
 \label{eq:tpc32-high-energy-slack}\\
 2(1-\beta)-\beta
 &=\frac{266}{400}-\frac{267}{400}
 =-\boxed{\frac1{400}}.
 \label{eq:tpc32-high-exception-density-exponent}
\end{align}
The $1/800$ quantity is an amplitude-exponent gap.  The $1/400$
quantity is its squared-energy counterpart and also the density exponent
for the exceptional frequency set at the large-content target scale.

\subsection{An unconditional density-one frequency theorem}

Take
\begin{equation}
 C\le J
 \label{eq:tpc32-high-C-below-J}
\end{equation}
and choose the no-wrap prime $q\asymp Q$ from
\cref{thm:tpc32-no-wrap-parseval}.  Then
\begin{equation}
 \left(
 \frac1q\sum_{r\, (q)}|\widehat A_{C,q}(r)|^2
 \right)^{1/2}
 \ll X^{o(1)}N_0X^{-267/800}.
 \label{eq:tpc32-high-rms-bound}
\end{equation}
Thus the rms determinant frequency has a saving larger than the
TPC-30 large-content margin by exactly $1/800$ in the exponent.

\begin{corollary}[High-$\beta$ almost-all-frequency saving]
\label{cor:tpc32-high-almost-all-frequencies}
Fix $0<\eta<133/400$ and put
\begin{equation}
 \sigma_\eta=\frac{133}{400}-\eta.
 \label{eq:tpc32-high-sigma-eta}
\end{equation}
For every fixed $\eps>0$,
\begin{align}
 &\#\left\{r\pmod q:
 |\widehat A_{C,q}(r)|>
 N_0X^{-133/400+\eta}\right\}
 \notag\\
 &\hspace{9em}
 \ll_{\eps,h,\mathscr D}
 X^{266/400-2\eta+\eps},
 \label{eq:tpc32-high-exception-count}\\
 &\frac1q
 \#\left\{r\pmod q:
 |\widehat A_{C,q}(r)|>
 N_0X^{-133/400+\eta}\right\}
 \notag\\
 &\hspace{9em}
 \ll_{\eps,h,\mathscr D}
 X^{-1/400-2\eta+\eps+o(1)}.
 \label{eq:tpc32-high-exception-proportion}
\end{align}
In particular, after choosing $\eps<\eta$, a proportion
$1-O(X^{-1/400-\eta+o(1)})$ of the frequencies has the stated saving.
The exceptional set may contain $r=0$.
\end{corollary}

\begin{proof}
Insert \eqref{eq:tpc32-high-sigma-eta} into
\eqref{eq:tpc32-no-wrap-exceptional-count}:
\[
 2\sigma_\eta
 =\frac{266}{400}-2\eta.
\]
This proves \eqref{eq:tpc32-high-exception-count}.  Since
$q=X^{267/400+o(1)}$, division by $q$ and
\eqref{eq:tpc32-high-exception-density-exponent} give
\eqref{eq:tpc32-high-exception-proportion}.
\end{proof}

The conclusion is stronger than merely saying that a generic frequency
has some power saving.  It says that the precise saving required to meet
the $133/400$ large-content floor already holds outside a set of relative
density $X^{-1/400+o(1)}$.  It does not identify the distinguished zero
frequency with the density-one set.

The Farey-family version has the same rms endpoint $\beta/2$.  Taking
$V=X^{267/800+o(1)}=Q^{1/2+o(1)}$ in
\cref{thm:tpc32-farey-large-sieve} gives, outside a power-small
proportion of reduced rational frequencies, every fixed saving below
\begin{equation}
 \min\left\{\frac{267}{800},\frac{267}{800}\right\}
 =\frac{267}{800}.
 \label{eq:tpc32-high-farey-rms-endpoint}
\end{equation}
Again, this family contains no copy of the zero rational frequency
because its denominators begin at two.

\subsection{The exact conditional threshold}

For the full-shell splice choose
\begin{equation}
 C=\lfloor J\rfloor=X^{133/400+o(1)},
 \qquad
 \kappa=\frac{133}{400}.
 \label{eq:tpc32-high-content-choice}
\end{equation}
The two large-content losses in
\eqref{eq:tpc32-conditional-full-shell-bound} then both have endpoint
$133/400$.  Because the conclusion asks for every strict saving below this
endpoint, the flatness term reaches it precisely when
\begin{align}
 \frac{\beta-\chi}{2}
 &\ge\frac{133}{400}
 \notag\\
 &\Longleftrightarrow
 \chi\le\beta-\frac{266}{400}
 =\frac{267}{400}-\frac{266}{400}
 =\boxed{\frac1{400}}.
 \label{eq:tpc32-high-chi-threshold}
\end{align}

\begin{corollary}[Conditional high-$\beta$ closure of the matched shell]
\label{cor:tpc32-high-conditional-closure}
Assume that the actual matched determinant coefficient with the choice
\eqref{eq:tpc32-high-content-choice} satisfies
\begin{equation}
 \mathfrak F_0(A_C)\le X^{\chi+o(1)}
 \qquad\text{for some fixed }\chi\le\frac1{400}.
 \label{eq:tpc32-high-flatness-input}
\end{equation}
Then every fixed
\begin{equation}
 \boxed{\sigma<\frac{133}{400}}
 \label{eq:tpc32-high-conditional-saving}
\end{equation}
is a power-saving exponent relative to $N_0\asymp XQ$ for the complete
matched raw cutoff shell.
\end{corollary}

\begin{proof}
Under \eqref{eq:tpc32-high-flatness-input},
\eqref{eq:tpc32-high-chi-threshold} shows that the flatness margin is at
least $133/400$.  The other two entries in
\eqref{eq:tpc32-conditional-saving-range} both equal $133/400$ by
\eqref{eq:tpc32-high-content-choice}.  Apply
\cref{thm:tpc32-conditional-zero-transfer} with a strict final
exponent.
\end{proof}

Condition \eqref{eq:tpc32-high-flatness-input} has a transparent
amplitude interpretation.  By definition,
\begin{equation}
 |\widehat A_{C,q}(0)|
 \le X^{\chi/2+o(1)}\|A_C\|_2.
 \label{eq:tpc32-high-amplitude-flatness}
\end{equation}
Thus the boundary $\chi\le1/400$ permits the zero coefficient to be larger
than the rms determinant coefficient by a factor at most
\begin{equation}
 X^{1/800+o(1)}.
 \label{eq:tpc32-high-amplitude-inflation}
\end{equation}
At $\chi=0$, the zero frequency is no larger than the rms scale up to
$X^{o(1)}$, and its saving endpoint is $267/800$; the final splice is
still limited to $266/800=133/400$ by large content.  By contrast,
the unconditional Cauchy bound permits
$\mathfrak F_0(A_C)\asymp Q=X^{267/400+o(1)}$, which gives no saving at
all.  Moving from that worst case to
\eqref{eq:tpc32-high-flatness-input} is the genuine arithmetic task;
the small numerical slack does not make it a soft consequence of the
frequency average.

\begin{center}
\begin{tabular}{>{\raggedright\arraybackslash}p{0.46\textwidth}
                >{\centering\arraybackslash}p{0.20\textwidth}
                >{\raggedright\arraybackslash}p{0.23\textwidth}}
\toprule
Quantity & Exact exponent & Status \\
\midrule
Large-content splice margin $1-\beta$
& $133/400=266/800$ & Unconditional \\
\addlinespace
Rms determinant-frequency margin $\beta/2$
& $267/800$ & Unconditional \\
\addlinespace
Rms amplitude slack over the splice
& $1/800$ & Exact arithmetic \\
\addlinespace
Bad-frequency density exponent at the splice target
& $1/400$ & Unconditional \\
\addlinespace
Allowed squared-flatness loss
& $\chi\le1/400$ & Additional input \\
\addlinespace
Complete matched-shell saving
& every $\sigma<133/400$ & Conditional on flatness \\
\bottomrule
\end{tabular}
\end{center}

The last row concerns only the matched raw cutoff shell delivered to
this paper.  Even under the flatness input it does not, by itself,
prove a prime-pair asymptotic, positivity, twin primes, or a breach of
the sieve parity barrier.
